% Do not chage this part
\begin{center}
  \textbf{{\Large ACKNOWLEDGEMENT}}\\[60pt]
\end{center}
\addcontentsline{toc}{chapter}{\textbf{\normalsize{\emph{ACKNOWLEDGEMENT}}}}
%%%%%%%%%%%%%%%%%%%%%%%%%%%%%%%%%%%%%%%

We wish to express our sincere gratitude to everyone who supported us throughout this undergraduate thesis.

First and foremost, we are deeply indebted to our thesis supervisor, \textbf{Dr. Md. Mostofa Akbar}, Professor, Department of Computer Science and Engineering, Bangladesh University of Engineering and Technology (BUET). His patient guidance, constructive criticism, and insistence on rigorous experimental methodology shaped every major design decision in VigiDroid---from the temporal evaluation protocol to the on-device metrics harness. Bi-weekly review meetings with him kept our research direction focused and academically sound.

We thank the faculty members and staff of the Department of Computer Science and Engineering, BUET, for providing the academic environment, laboratory facilities, and administrative support that made this work possible. The undergraduate curriculum gave us the foundations in algorithms, machine learning, and software engineering upon which this project was built; the thesis experience taught us how to extend those foundations into an integrated research system.

We acknowledge the maintainers of the public datasets and tools that underpin our evaluation: the AndroZoo corpus, the Hugging Face \texttt{android-apks} dataset, the authors of the ONNX Runtime mobile stack, and the open-source communities behind PyTorch, scikit-learn, and XGBoost. Their work enabled reproducible research at a scale that would otherwise have been impractical for an undergraduate group.

We are grateful to our families for their encouragement during the long hours spent debugging parity failures, staging ONNX bundles, and running overnight device evaluation batches. Their patience and belief in us sustained the project through its most difficult phases.

Finally, we thank one another---Abdullah Faiyaz, Jarin Tasnim, and Mohammad Sadnam Faiyaz---for the trust, division of labour, and honest peer review that transformed an ambitious idea into a working on-device malware detection framework.

\vspace*{20.0mm}

\begin{minipage}[t]{0.2\textwidth}
  Dhaka\par
  \thesisdate
\end{minipage}%
\hfill
\begin{minipage}[t]{0.45\textwidth}
  \begin{enumerate}
    \vspace{-0.75\baselineskip}
    \DTLforeach{ThesisStudents}{\StudentName=Column1}{\item[]\StudentName}
  \end{enumerate}
\end{minipage}

\endinput
